\documentclass[11pt,a4paper]{article}
\usepackage[utf8]{inputenc}\usepackage[T1]{fontenc}\usepackage[italian,english]{babel}
\usepackage{amsmath,amssymb,amsthm}\usepackage[margin=2.2cm]{geometry}\usepackage{longtable,booktabs,array,calc}
\usepackage{listings,xcolor}\usepackage{hyperref}\usepackage{url}
\providecommand{\tightlist}{\setlength{\itemsep}{0pt}\setlength{\parskip}{0pt}}
\lstset{basicstyle=\ttfamily\scriptsize,breaklines=true,frame=single,captionpos=b,columns=fullflexible,keepspaces=true,inputencoding=utf8,extendedchars=true}
\title{Proof Pursuit --- Problem 1, part 6\\ \large prove, decisioni degli agenti, codice verificato, letteratura}
\author{Team Proof Pursuit (Researcher: T. Tumini; Referee: gabundos; Creative: M. Renzi; gio3r3gio)}
\date{\today}
\begin{document}\maketitle\tableofcontents
\section*{Come leggere questo rapporto}
Per ogni cella: la richiesta ufficiale, lo stato dichiarato dal team (mai ``risolto'' senza revisione umana), la consegna
con la prova, poi la traccia delle decisioni degli agenti (Researcher: famiglia e motivo; Referee: verdetto ed errore
fatale; Creative: quando e perché è intervenuto) e il codice su cui il tentativo poggia con l'esito della riesecuzione
fidata. DIMOSTRATO, CITATO, VERIFICATO SU INTERVALLO FINITO e CONGETTURA restano distinti. L'architettura del sistema
è descritta nel README del repository.
\section{Problem 1 — Angles between lines (Fejes Toth)}
\subsection{Cella 6 (13 punti) — stato: parziale}
\subparagraph{\texorpdfstring{Parte 6 --- \(N\) lines in
\(\mathbb{R}^d\)}{Parte 6 --- N lines in \textbackslash mathbb\{R\}\^{}d}}\label{parte-6-n-lines-in-mathbbrd}

\textbf{Punteggio:} 13 points · \textbf{Valutazione:} Judged ·
\textbf{Open question}

Fejes Tóth's conjecture from the Setting, for every \(N\) and every
\(d\). For \(N\) lines in \(\mathbb{R}^d\), write \(N = qd + s\) with
\(0 \le s < d\), and let \[
M(N,d) = s \binom{q+1}{2} + (d-s) \binom{q}{2}.
\] Prove or disprove: every \(N\) lines in \(\mathbb{R}^d\) satisfy \[
S \;\le\; \left( \binom{N}{2} - M(N,d) \right) \frac{\pi}{2}.
\] This is the value attained by splitting the lines as evenly as
possible among \(d\) mutually orthogonal directions. Settling any
infinite family not already covered above counts as partial progress.

\emph{(Fine del problema 1. Punteggi: 1+2+3+5+8+13 = 32.)}

\textbf{Enunciato protetto dato agli agenti (state.json).} \texttt{OPEN. For N lines in R\textasciicircum{}d write N = qd+s, 0 \textless{}= s \textless{} d, M(N,d) = s C(q+1,2) + (d-s) C(q,2). Prove or disprove: S \textless{}= (C(N,2) - M(N,d)) pi/2. Settling any infinite family not covered by cells 1-5 counts as partial progress.}

\textbf{Evidenza registrata in STATUS.md.} loop p1\_c6 attempt\_001: REJECT (The only declared claim is `main` (Fejes Tóth's bound for all N, d). The submission explicitly states in §8 that no non-); lemmi intermedi in submission/parte-6.en.md

\subsubsection{Consegna (prova)}
\paragraph{Problem 1 --- Part 6 --- submission
(PARTIAL)}\label{problem-1-part-6-submission-partial}

\textbf{Declared status: PARTIAL.} The cell is not solved. What follows
is the intermediate progress actually established, as judged by the
automatic Referee (verdict \texttt{REJECT}: The only declared claim is
\texttt{main} (Fejes Tóth's bound for all N, d). The submission
explicitly states in §8 that no non-trivial infinite family is proved
and that cell 6 remains open; Lemmas 0--4 and Corollaries 5--6 are
conditional transfer statements (FT(N,d) ⇒ FT(N+1,d) for s=d−1;
FT(N+1,d) ⇒ FT(N,d) for d \textbar{} N; FT(N+q',d) ⇒ FT(N,d−1)) whose
hypotheses are themselves the unproved conjecture. None ). The author's
own declared status was \texttt{LEMMA\_CANDIDATE}.

\subparagraph{1. Result and scope}\label{result-and-scope}

Rigorous, self-contained proofs of: (i) exact identities for M(N,d);
(ii) the trivial infinite family N ≤ d (and d = 1), not covered by cells
1--5 but of no depth; (iii) three transfer lemmas: FT(N,d) ⟹ FT(N+1,d)
when N ≡ −1 (mod d); FT(N+1,d) ⟹ FT(N,d) when d \textbar{} N; FT(N+q,d)
⟹ FT(N,d−1) with q = ⌊N/(d−1)⌋ (finite-N analogue of Bilyk--Matzke's
continuous dimension reduction); (iv) conditional corollary: cell 4's
(5,3) implies S ≤ 6π for six lines in R\^{}3. Float search in
R\textsuperscript{3..R}6 (13 pairs) found no counterexample. No
non-trivial infinite family is settled; cell 6 remains open.

\subparagraph{2. Proof}\label{proof}

\paragraph{Cell 6 --- transfer lemmas, the trivial family, and a
counterexample
search}\label{cell-6-transfer-lemmas-the-trivial-family-and-a-counterexample-search}

\subparagraph{0. Notation}\label{notation}

Lines are lines through the origin of \(\mathbb R^d\), spanned by unit
vectors;
\(\theta(\ell,\ell')=\arccos|\langle x,x'\rangle|\in[0,\pi/2]\);
\(S=\sum_{i<j}\theta(\ell_i,\ell_j)\). For integers \(N\ge1\), \(d\ge1\)
write \(N=qd+s\) with \(q=\lfloor N/d\rfloor\), \(0\le s<d\), and
\[M(N,d)=s\binom{q+1}{2}+(d-s)\binom{q}{2},\qquad B(N,d)=\Big(\binom N2-M(N,d)\Big)\frac{\pi}{2}.\]
\textbf{Definition.} \(\mathrm{FT}(N,d)\) is the statement: \emph{every
\(N\) lines in \(\mathbb R^d\) satisfy \(S\le B(N,d)\).} Cell 6 asks
whether \(\mathrm{FT}(N,d)\) holds for all \(N,d\).

Throughout, \(\binom{0}{2}=\binom{1}{2}=0\).

\subparagraph{\texorpdfstring{1. Exact identities for
\(M\)}{1. Exact identities for M}}\label{exact-identities-for-m}

\textbf{Lemma 0.} Let \(N=qd+s\), \(0\le s<d\). (a)
\(2M(N,d)=q\,(N+s-d)\). (b) \(M(N+1,d)=M(N,d)+q\). (c) If \(d\ge2\) and
\(N=q'(d-1)+s'\) with \(0\le s'<d-1\) (so
\(q'=\lfloor N/(d-1)\rfloor\)), then
\(M(N+q',d)=M(N,d-1)+\binom{q'}{2}\).

\emph{Proof.} (a)
\(M=s\frac{q(q+1)}2+(d-s)\frac{q(q-1)}2=\frac q2\big[s(q+1)+(d-s)(q-1)\big]=\frac q2\big[dq+2s-d\big]=\frac q2(N+s-d)\),
using \(dq+s=N\).

\begin{enumerate}
\def\labelenumi{(\alph{enumi})}
\setcounter{enumi}{1}
\item
  If \(s\le d-2\), then \(N+1=qd+(s+1)\) with \(0\le s+1<d\) is the
  canonical representation, and
  \(M(N+1,d)-M(N,d)=\binom{q+1}2-\binom q2=q\). If \(s=d-1\), then
  \(N+1=(q+1)d+0\), so \(M(N+1,d)=d\binom{q+1}2\), while
  \(M(N,d)=(d-1)\binom{q+1}2+\binom q2\); the difference is again
  \(\binom{q+1}2-\binom q2=q\).
\item
  \(N+q'=q'(d-1)+s'+q'=q'd+s'\) with \(0\le s'<d-1<d\), so this is the
  canonical representation of \(N+q'\) modulo \(d\), and
  \(M(N+q',d)=s'\binom{q'+1}2+(d-s')\binom{q'}2=\Big[s'\binom{q'+1}2+(d-1-s')\binom{q'}2\Big]+\binom{q'}2=M(N,d-1)+\binom{q'}2.\)
  \(\square\)
\end{enumerate}

(These identities were also checked by exact integer arithmetic for all
\(1\le d\le40\), \(1\le N\le400\) --- file
\texttt{verifica\_identita\_M.py}; the proof above does not depend on
that check.)

\subparagraph{\texorpdfstring{2. The trivial infinite family \(N\le d\)
(and
\(d=1\))}{2. The trivial infinite family N\textbackslash le d (and d=1)}}\label{the-trivial-infinite-family-nle-d-and-d1}

\textbf{Lemma 1.} \(\mathrm{FT}(N,d)\) holds whenever \(N\le d\), and
whenever \(d=1\).

\emph{Proof.} If \(N\le d\) then \(q=0\), \(s=N\), so
\(M(N,d)=N\binom12+(d-N)\binom02=0\) and \(B(N,d)=\binom N2\frac\pi2\).
Since every \(\theta\le\pi/2\), \(S\le\binom N2\frac\pi2\). If \(d=1\),
all lines equal \(\mathbb R^1\), so \(S=0\); here \(N=N\cdot1+0\),
\(M(N,1)=\binom N2\), \(B(N,1)=0\), and \(S\le0\) holds. \(\square\)

\emph{Honest remark.} This family is not covered by cells 1--5 (cell 3
starts at \(N=d+1\)) but it is trivial; I list it only for completeness,
not as substantive progress.

\subparagraph{\texorpdfstring{3. Adding a line when
\(N\equiv-1\pmod d\)}{3. Adding a line when N\textbackslash equiv-1\textbackslash pmod d}}\label{adding-a-line-when-nequiv-1pmod-d}

\textbf{Lemma 2.} Let \(d\ge1\), \(N\ge2\), \(N=qd+s\) with \(s=d-1\)
(i.e.~\(N\equiv-1\pmod d\)). If \(\mathrm{FT}(N,d)\) holds, then
\(\mathrm{FT}(N+1,d)\) holds.

\emph{Proof.} Let \(\ell_1,\dots,\ell_{N+1}\) be lines in
\(\mathbb R^d\) and \(S\) their angle sum. For \(k=1,\dots,N+1\) let
\(S_k\) be the angle sum of the \(N\) lines obtained by deleting
\(\ell_k\). A pair \(\{i,j\}\) contributes \(\theta(\ell_i,\ell_j)\) to
\(S_k\) exactly for the \(N+1-2=N-1\) indices \(k\notin\{i,j\}\), hence
\[\sum_{k=1}^{N+1}S_k=(N-1)\,S.\] By \(\mathrm{FT}(N,d)\),
\(S_k\le B(N,d)\) for each \(k\), so \((N-1)S\le(N+1)B(N,d)\),
i.e.~\(S\le\frac{N+1}{N-1}B(N,d)\) (here \(N-1\ge1\)). It remains to
show \(\frac{N+1}{N-1}B(N,d)\le B(N+1,d)\), i.e.
\[(N+1)\Big(\binom N2-M(N,d)\Big)\le(N-1)\Big(\binom{N+1}2-M(N+1,d)\Big).\tag{1}\]
Since \((N+1)\binom N2=\frac{(N+1)N(N-1)}2=(N-1)\binom{N+1}2\), and
\(M(N+1,d)=M(N,d)+q\) by Lemma 0(b), (1) is equivalent to
\(-(N+1)M(N,d)\le-(N-1)M(N,d)-(N-1)q\), i.e.~to \((N-1)q\le2M(N,d)\). By
Lemma 0(a), \(2M(N,d)=q(N+s-d)=q(N-1)\) because \(s=d-1\). So (1) holds
with equality, and \(S\le B(N+1,d)\). \(\square\)

\emph{Remark.} If \(s\ne d-1\) and \(q\ge1\), the inequality
\((N-1)q\le q(N+s-d)\) fails, so this averaging argument gives exactly
the residue class \(s=d-1\) and nothing else (exact check in
\texttt{verifica\_identita\_M.py}, item (d)).

\subparagraph{\texorpdfstring{4. Removing a line when
\(d\mid N\)}{4. Removing a line when d\textbackslash mid N}}\label{removing-a-line-when-dmid-n}

\textbf{Lemma 3.} Let \(d\ge1\), \(N=qd\) with \(q\ge1\). If
\(\mathrm{FT}(N+1,d)\) holds, then \(\mathrm{FT}(N,d)\) holds.

\emph{Proof.} If \(N=1\) then \(S=0\le B(1,d)\). Let \(N\ge2\) and let
\(\ell_1,\dots,\ell_N\) be lines in \(\mathbb R^d\) with angle sum
\(S\). For each \(k\) put \(A_k=\sum_{j\ne k}\theta(\ell_k,\ell_j)\).
Each unordered pair \(\{i,j\}\) appears in exactly \(A_i\) and \(A_j\),
so \(\sum_kA_k=2S\) and some index \(k\) satisfies \(A_k\ge2S/N\).
Consider the \(N+1\) lines \(\ell_1,\dots,\ell_N,\ell_{N+1}\) with
\(\ell_{N+1}=\ell_k\) (a repeated line, allowed by the problem). Its
angle sum is
\(S'=S+\sum_{j\le N}\theta(\ell_{N+1},\ell_j)=S+A_k+\theta(\ell_k,\ell_k)=S+A_k\ge S\,(1+\tfrac2N)=\tfrac{N+2}{N}S\),
since \(\theta(\ell_k,\ell_k)=\arccos1=0\). By \(\mathrm{FT}(N+1,d)\),
\(S'\le B(N+1,d)\), hence \(S\le\frac{N}{N+2}B(N+1,d)\). It remains to
show \(\frac N{N+2}B(N+1,d)\le B(N,d)\), i.e.
\[N\Big(\binom{N+1}2-M(N+1,d)\Big)\le(N+2)\Big(\binom N2-M(N,d)\Big).\tag{2}\]
Compute
\(N\binom{N+1}2-(N+2)\binom N2=\frac{N}{2}\big[(N^2+N)-(N^2+N-2)\big]=N\).
With Lemma 0(b), (2) is equivalent to
\(N\le N\,(M(N,d)+q)-(N+2)M(N,d)=Nq-2M(N,d)\), and by Lemma 0(a) with
\(s=0\): \(Nq-2M(N,d)=Nq-q(N-d)=qd=N\). So (2) holds with equality and
\(S\le B(N,d)\). \(\square\)

\emph{Remark.} For \(s\ne0\) one gets \(Nq-2M=q(d-s)<qd+s=N\), so the
argument works exactly when \(d\mid N\) (exact check, item (e)).

\subparagraph{5. Dropping one dimension}\label{dropping-one-dimension}

\textbf{Lemma 4.} Let \(d\ge2\), \(N\ge1\),
\(q'=\lfloor N/(d-1)\rfloor\). If \(\mathrm{FT}(N+q',d)\) holds, then
\(\mathrm{FT}(N,d-1)\) holds.

\emph{Proof.} Let \(\ell_1,\dots,\ell_N\) be lines in
\(\mathbb R^{d-1}\) spanned by unit vectors \(x_1,\dots,x_N\), with
angle sum \(S\). Embed \(\iota:\mathbb R^{d-1}\to\mathbb R^d\),
\(\iota(x)=(x,0)\); this is a linear isometry, so
\(\langle\iota x_i,\iota x_j\rangle=\langle x_i,x_j\rangle\) and all
angles \(\theta\) are preserved. Let \(e_d=(0,\dots,0,1)\) and consider
the \(N+q'\) lines in \(\mathbb R^d\) spanned by
\(\iota x_1,\dots,\iota x_N\) and \(q'\) copies of \(e_d\). Their angle
sum is
\[S'=S+\sum_{i=1}^{N}q'\,\theta(\iota x_i,e_d)+\binom{q'}2\theta(e_d,e_d)=S+q'N\frac\pi2+0,\]
because \(\langle\iota x_i,e_d\rangle=0\) gives
\(\theta=\arccos0=\pi/2\) and \(\theta(e_d,e_d)=0\). By
\(\mathrm{FT}(N+q',d)\),
\[S\le\Big(\binom{N+q'}2-M(N+q',d)-q'N\Big)\frac\pi2.\] Now
\(\binom{N+q'}2-q'N=\frac{(N+q')(N+q'-1)}2-q'N=\frac{N^2-N+q'^2-q'}2=\binom N2+\binom{q'}2\),
and \(M(N+q',d)=M(N,d-1)+\binom{q'}2\) by Lemma 0(c). Therefore
\(S\le\big(\binom N2-M(N,d-1)\big)\frac\pi2=B(N,d-1)\). \(\square\)

\subparagraph{6. Consequences}\label{consequences}

\textbf{Corollary 5 (conditional).} If cell 4's statement for five lines
in \(\mathbb R^3\) holds (\(\mathrm{FT}(5,3)\), i.e.~\(S\le4\pi\)), then
every six lines in \(\mathbb R^3\) satisfy \(S\le6\pi\),
i.e.~\(\mathrm{FT}(6,3)\). \emph{Proof.} \(N=5=1\cdot3+2\), so
\(s=2=d-1\) and \(N\ge2\); Lemma 2 gives \(\mathrm{FT}(6,3)\). Here
\(B(6,3)=(15-3)\frac\pi2=6\pi\) since \(6=2\cdot3+0\) and
\(M(6,3)=3\binom22=3\). \(\square\) (This recovers, conditionally on
cell 4, the range \(N\le6\) in \(\mathbb R^3\) attributed by
Bilyk--Matzke to Fejes Tóth 1959; cell 4 itself is not proved here.)

\textbf{Corollary 6.} In the plane, Lemma 2 with \(d=2\) says
\(\mathrm{FT}(N,2)\Rightarrow\mathrm{FT}(N+1,2)\) for odd \(N\ge3\), and
Lemma 3 says \(\mathrm{FT}(N+1,2)\Rightarrow\mathrm{FT}(N,2)\) for even
\(N\). (Consistent with cell 1, which is already proved by the team; not
needed.)

\textbf{What the lemmas do not give.} Starting from the families of
cells 3--5 (\(N=d+1\), \(N=d+2\)), Lemma 2 applies only when
\(d+2\equiv-1\pmod d\), i.e.~\(d\mid3\), and Lemma 4 maps \((d+2,d)\) to
\((d+1,d-1)\), the same family. So no new \emph{infinite} family follows
from cells 1--5 plus these lemmas; they are transfer tools for future
attempts (e.g.~any proof of \(\mathrm{FT}(2d-1,d)\) would yield
\(\mathrm{FT}(2d,d)\) by Lemma 2, and \(\mathrm{FT}(2d,d)\) for all
\(d\) would yield \(\mathrm{FT}(2d-2,d-1)\) by Lemma 4).

\subparagraph{7. Counterexample search (float, exploration only ---
proves
nothing)}\label{counterexample-search-float-exploration-only-proves-nothing}

Using the team's hill-climbing harness (\texttt{tools/autoloop.py} on
\texttt{problema-1/esperimenti/lines\_Rd.py}; Gaussian perturbations
plus structured moves: copy a vector, orthogonalise against a random
subset, random jump; 3 seeds × 8 restarts × 5 s per pair), the best
\(S\) found versus \(B(N,d)\) for
\((d,N)\in\{(3,6),(3,7),(3,8),(3,9),(4,7),(4,8),(4,9),(4,10),(5,8),(5,9),(5,11),(6,10),(6,13)\}\)
(each \(\approx5\) s wall, \(\approx2.3\)--\(3.2\cdot10^5\) iterations):
- in 12 of 13 cases the best value agrees with \(B(N,d)\) to within
\(2\cdot10^{-7}\); the sign of the residual is negative
(\(\approx-10^{-7}\)) in these cases, which is the expected
floating-point artefact of \texttt{arccos} at arguments clipped to
\(1-\varepsilon\)
(\(\arccos(1-\varepsilon)\approx\sqrt{2\varepsilon}\)), \textbf{not}
evidence of a violation; - for \((3,6)\) the search did not converge
(best \(18.8287<6\pi=18.8496\)). No configuration exceeding the
conjectured bound was found. This is float arithmetic on a finite random
sample: it supports ``prove'' over ``disprove'' heuristically and
settles nothing.

\subparagraph{8. Summary of status}\label{summary-of-status}

\begin{itemize}
\tightlist
\item
  Proved (rigorously, self-contained): Lemma 0 (identities), Lemma 1
  (trivial family \(N\le d\) and \(d=1\)), Lemmas 2--4 (transfer
  implications), Corollaries 5--6.
\item
  Not proved: any non-trivial infinite family; cell 6 remains open.
  Nothing here is claimed as a solution of the cell.
\end{itemize}

\subparagraph{3. Verification: instructions, dependencies,
timings}\label{verification-instructions-dependencies-timings}

Python 3 standard library. Scripts (re-run by the orchestrator, see
\texttt{runs/p1\_c6/verifica/}): - code\_1 (python, rigor
\texttt{exact}): Exact integer check of the identities of Lemma 0 and of
the equality/inequality conditions in Lemmas 2--3, for 1≤d≤40, 1≤N≤400
(0.1 s). Saved as runs/p1\_c6/sandbox/verifica\_identita\_M.py. The
proofs do not depend on it. - code\_2 (bash, rigor
\texttt{float\_exploration\_only}): Float-only hill-climbing search for
configurations exceeding the conjectured bound in
R\textsuperscript{3..R}6 (13 pairs (d,N), 3 seeds x 8 restarts x 5 s
each, \textasciitilde70 s total), using the team's existing harness
tools/autoloop.py with objective problema-1/esperimenti/lines\_Rd.py;
ledger saved in runs/p1\_c6/sandbox/esplora\_Rd.jsonl. Proves nothing.

\subparagraph{4. Sources and
contribution}\label{sources-and-contribution}

\begin{itemize}
\tightlist
\item
  arXiv:1801.07837 (Bilyk--Matzke) --- context only: Prop. 3.1
  (dimension reduction for the continuous energy) motivates Lemma 4; the
  `only settled case is d=1' statement; nothing cited as a proof step
\item
  problema-1/fonti/README.md (team reading notes on Bilyk--Matzke;
  attribution of N≤6 in R\^{}3 to Fejes Tóth 1959, unread original)
\item
  Position with respect to the literature: {[}1801.07837{]}
  Bilyk--Matzke (read in full by the team earlier, see
  problema-1/fonti/README.md): the only settled case is the plane; they
  give a general non-sharp energy bound π/2 − 69/(50(d+1)) via the frame
  potential and, for the continuous (measure) version, a
  dimension-reduction Prop. 3.1 (conjecture in S\^{}d ⟹ conjecture in
  S\^{}\{d−1\}). Lemma 4 below is the finite-N analogue of that
  reduction, proved from scratch by embedding and adding q copies of a
  new orthogonal axis; the exact identity M(N+q,d) = M(N,d−1) + C(q,2)
  is what makes it sharp for finite N. Lemmas 2--3 (averaging over
  deleting/duplicating a line) do not appear in the abstracts and are
  our own. {[}2007.08698{]} Lim--McCann (abstract only): deformation to
  α-powers, optimality for α = ∞; not used. {[}1204.3850{]} is
  irrelevant (polygon mapping). Nothing in the literature I have settles
  any new infinite family; my approach departs from the energy method
  (known non-sharp) and instead builds exact transfer lemmas.
\end{itemize}

\subparagraph{5. Limits and unresolved
parts}\label{limits-and-unresolved-parts}

Referee's blocking point: The only declared claim is \texttt{main}
(Fejes Tóth's bound for all N, d). The submission explicitly states in
§8 that no non-trivial infinite family is proved and that cell 6 remains
open; Lemmas 0--4 and Corollaries 5--6 are conditional transfer
statements (FT(N,d) ⇒ FT(N+1,d) for s=d−1; FT(N+1,d) ⇒ FT(N,d) for d
\textbar{} N; FT(N+q',d) ⇒ FT(N,d−1)) whose hypotheses are themselves
the unproved conjecture. None of them, alone or combined with cells 1--5
(which are moreover not verified in the state: highest\_verified\_cell =
0), yields FT(N,d) for any family beyond the trivial N ≤ d / d = 1 case.
Hence the declared target is not established, and no other declared
claim exists to accept. The float search of §7 is exploration only and
proves nothing (correctly acknowledged by the author). Next step
required: Address the stated blocking obligation without silently
changing the target Gaps declared by the author: - No non-trivial
infinite family is settled; the only `family' proved outright (N ≤ d, d
= 1) is trivial and should not be scored as substantive progress on cell
6. - Corollary 5 (six lines in R\^{}3) is conditional on cell 4 (five
lines in R\^{}3), which is not proved in this attempt and is not among
the verified claims. - Section 7 is floating-point hill climbing on a
finite random sample: it is exploration only and proves nothing; the
(3,6) run did not even converge to the conjectured value. - The transfer
lemmas, combined with cells 1--5, do not yield any new infinite family
(shown in §6); their value is as reusable tools for later attempts.

\subsubsection{Decisioni degli agenti}
\paragraph{Run \texttt{p1\_c6}}
\begin{description}
\item[Iterazione 1 — attempt\_001]
\textbf{Researcher}: famiglia \texttt{reduction}, sotto-obiettivo: Establish rigorous, reusable reduction lemmas between instances FT(N,d) of the conjecture (add a line, remove a line, drop one dimension), settle the (trivial) infinite family N ≤ d not covered by cells 1–5, and record a float-only counterexample search in R\textasciicircum{}3..R\textasciicircum{}6.. Stato dichiarato: \texttt{LEMMA\_CANDIDATE}.
\emph{Perché questa scelta}: No blocker and no verified claims exist for this cell; cell 6 is an open problem and asks for infinite families beyond cells 1–5. A frontal attack (e.g. N = 2d) is out of reach in one iteration, and the frame-potential/Gram approach is known (Bilyk–Matzke) to be non-sharp. The best first step is to build the exact 'transfer' machinery: which instances of the conjecture imply which others. These lemmas are cheap, fully rigorous, reusable by later iterations (any future family proved for N ≡ −1 mod d extends to N+1; any family in R\textasciicircum{}d transfers down a dimension), and they immediately give a new finite case ((6,3) from cell 4). I also record honestly the trivial family N ≤ d, and a float-only search giving no counterexample, since the cell says 'prove or disprove'.
\emph{Rispetto allo stato dell'arte}: [1801.07837] Bilyk–Matzke (read in full by the team earlier, see problema-1/fonti/README.md): the only settled case is the plane; they give a general non-sharp energy bound π/2 − 69/(50(d+1)) via the frame potential and, for the continuous (measure) version, a dimension-reduction Prop. 3.1 (conjecture in S\textasciicircum{}d ⟹ conjecture in S\textasciicircum{}\{d−1\}). Lemma 4 below is the finite-N analogue of that reduction, proved from scratch by embedding and adding q copies of a new orthogonal axis; the exact identity M(N+q,d) = M(N,d−1) + C(q,2) is what makes it sharp for finite N. Lemmas 2–3 (averaging over deleting/duplicating a line) do not appear in the abstracts and are our own. [2007.08698] Lim–McCann (abstract only): deformation to α-powers, optimality for α = ∞; not used. [1204.3850] is irrelevant (polygon mapping). Nothing in the literature I have settles any new infinite family; my approach departs from the energy method (known non-sharp) and instead builds exact transfer lemmas.
\textbf{Referee}: verdetto \texttt{REJECT} (review\_status \texttt{NEEDS\_WORK}).
\emph{Errore fatale}: The only declared claim is `main` (Fejes Tóth's bound for all N, d). The submission explicitly states in §8 that no non-trivial infinite family is proved and that cell 6 remains open; Lemmas 0–4 and Corollaries 5–6 are conditional transfer statements (FT(N,d) ⇒ FT(N+1,d) for s=d−1; FT(N+1,d) ⇒ FT(N,d) for d | N; FT(N+q',d) ⇒ FT(N,d−1)) whose hypotheses are themselves the unproved conjecture. None of them, alone or combined with cells 1–5 (which are moreover not verified in the state: highest\_verified\_cell = 0), yields FT(N,d) for any family beyond the trivial N ≤ d / d = 1 case. Hence the declared target is not established, and no other declared claim exists to accept. The float search of §7 is exploration only and proves nothing (correctly acknowledged by the author).
\emph{Prossimo passo richiesto}: Address the stated blocking obligation without silently changing the target
\end{description}

\subsubsection{Codice su cui poggia il tentativo}
\begin{lstlisting}[caption={attempt\_001 code\_1 [python, rigore exact] — Exact integer check of the identities of Lemma 0 and of the equality/inequality conditions — verdetto del Referee: REJECT}]
"""
Verifica ESATTA (interi Python) delle identita' aritmetiche usate nei lemmi di riduzione
per la congettura di Fejes Toth (problema 1, parte 6).
  M(N,d) = s*C(q+1,2) + (d-s)*C(q,2),  N = q*d + s, 0 <= s < d.
Cosa controlla e perche':
  (a) 2*M(N,d) = q*(N+s-d)                      -> usata nei lemmi 2 e 3
  (b) M(N+1,d) = M(N,d) + q                     -> usata nei lemmi 2 e 3
  (c) M(N+q',d) = M(N,d-1) + C(q',2), q'=N//(d-1) -> usata nel lemma 4
  (d) disuguaglianza del lemma 2 vale sse s = d-1 oppure q = 0 (N>=2)
  (e) disuguaglianza del lemma 3 vale sse s = 0    (N>=1)
Copertura: 1 <= d <= 40, 1 <= N <= 400. Aritmetica esatta: nessun float.
"""
from math import comb


def M(N, d):
    """Numero di coppie coincidenti nella configurazione equidistribuita."""
    q, s = divmod(N, d)
    return s * comb(q + 1, 2) + (d - s) * comb(q, 2)


def B2(N, d):
    """Bound congetturato in unita' di pi/2: C(N,2) - M(N,d)."""
    return comb(N, 2) - M(N, d)


def controlla(d_max=40, N_max=400):
    for d in range(1, d_max + 1):
        for N in range(1, N_max + 1):
            q, s = divmod(N, d)
            assert 2 * M(N, d) == q * (N + s - d), (N, d)
            assert M(N + 1, d) == M(N, d) + q, (N, d)
            if d >= 2:
                qq = N // (d - 1)
                assert M(N + qq, d) == M(N, d - 1) + comb(qq, 2), (N, d)
            if N >= 2:
                lemma2 = (N + 1) * B2(N, d) <= (N - 1) * B2(N + 1, d)
                assert lemma2 == (s == d - 1 or q == 0), (N, d)
            lemma3 = N * B2(N + 1, d) <= (N + 2) * B2(N, d)
            assert lemma3 == (s == 0), (N, d)
    print(f"OK: identita' (a)-(e) verificate per d<=%d, N<=%d" % (d_max, N_max))


if __name__ == "__main__":
    controlla()

\end{lstlisting}
\noindent\emph{Riesecuzione fidata dell'orchestratore}: \texttt{orchestrator re-ran code\_1.py (python3, clean copy of the researcher sandbox): exit 0 in 0.1s; stdout: "OK: identita' (a)-(e) verificate per d\textless{}=40, N\textless{}=400"; stderr: ''}
\begin{lstlisting}[caption={attempt\_001 code\_2 [bash, rigore float\_exploration\_only] — Float-only hill-climbing search for configurations exceeding the conjectured bound in R\textasciicircum{}3. — verdetto del Referee: REJECT}]
cd runs/p1_c6/sandbox
for dN in "3 6" "3 7" "3 8" "3 9" "4 7" "4 8" "4 9" "4 10" "5 8" "5 9" "5 11" "6 10" "6 13"; do
  d=${dN%% *}; N=${dN##* }
  ../../../.venv/bin/python ../../../tools/autoloop.py ../../../problema-1/esperimenti/lines_Rd.py \
      --budget 5 --seeds 3 --restarts 8 --ledger esplora_Rd.jsonl --tag "c6_d${d}_N${N}" -- $d $N
done
# riassunto: per ogni (d,N) best trovato, target congetturato, gap (float)
../../../.venv/bin/python - <<'EOF'
import json
for riga in open("esplora_Rd.jsonl"):
    r = json.loads(riga)
    print(f"d={r['args'][0]} N={r['args'][1]:>2}  best={r['best_score']:.6f}  target={r['target']:.6f}  target-best={r['gap_target_minus_best']:+.2e}")
EOF
# Output osservato: (3,6) best 18.828661 < 18.849556 (non convergente); tutti gli altri 12 casi best = target entro 2e-7
# (residuo negativo ~1e-7 = artefatto float di arccos vicino a 1, non una violazione).
\end{lstlisting}
\noindent\emph{Riesecuzione fidata dell'orchestratore}: \texttt{code\_2: not re-run (language 'bash' not supported by the orchestrator)}

\section{arXiv literature consulted}
\begin{thebibliography}{99}
\bibitem{1801.07837v1} Dmitriy Bilyk, Ryan W Matzke. \emph{On the Fejes Tóth Problem about the Sum of Angles Between Lines}. arXiv:1801.07837v1 (2018). \url{http://arxiv.org/abs/1801.07837v1}
\bibitem{2007.08698v2} Tongseok Lim, Robert J. McCann. \emph{On Fejes Tóth's conjectured maximizer for the sum of angles between lines}. arXiv:2007.08698v2 (2020). \url{http://arxiv.org/abs/2007.08698v2}
\bibitem{1204.3850v1} Jérémie Chalopin, Shantanu Das, Yann Disser, Matúš Mihalák, Peter Widmayer. \emph{Simple Agents Learn to Find Their Way: An Introduction on Mapping Polygons}. arXiv:1204.3850v1 (2012). \url{http://arxiv.org/abs/1204.3850v1}
\end{thebibliography}
\end{document}
